\documentclass[11pt,a4paper]{article}
\usepackage[T1]{fontenc}
\usepackage{booktabs}
\usepackage[colorlinks=true,urlcolor=blue,linkcolor=blue,citecolor=blue]{hyperref}
\setlength{\oddsidemargin}{-0.24cm}
\setlength{\evensidemargin}{-0.24cm}
\setlength{\topmargin}{-0.24cm}
\setlength{\headheight}{0pt}
\setlength{\headsep}{0pt}
\setlength{\textwidth}{16.4cm}
\setlength{\textheight}{25.1cm}
\setlength{\parindent}{0pt}
\setlength{\parskip}{0.38em}
\newcommand{\mfl}{m_{4\ell}}
\newcommand{\W}{W_{68}}
\newcommand{\Pois}{\mathrm{Pois}}
\newcommand{\arxivref}[1]{\href{https://arxiv.org/abs/#1}{arXiv:#1}}
\begin{document}
\begin{center}
{\Large\bfseries Paper outline for supervisory discussion}\\[0.4em]
{\large Kinematic feature groups and the robustness of signal-strength inference}\\
in simulated $H\to ZZ^*\to2e2\mu$ events \hfill 5 October 2026
\end{center}

\textbf{Summary.} This proposed study asks how kinematic input groups affect
expected signal-strength interval widths within a fixed learning and inference
procedure, and how robust the differences are. The exploratory comparison
identifies compact candidates and different leading configurations under
classification AUC and nominal interval width. Their advantage over the full
input remains inconclusive under finite-MC resampling, while coverage needs
further validation. Two priorities for discussion are the suitability of the
educational data release for a publication study and the evidence needed to
support the paper's main conclusion.

\section{Research question and intended contribution}
The primary question is how input groups change the expected 68\% interval
width for signal strength $\mu$, which scales the simulated signal yield,
and whether those changes withstand training randomness, finite Monte Carlo
(MC) samples and alternative inference settings. Only the first two sources
have existing diagnostics; robustness to alternative inference settings remains
a research goal. AUC comparisons and compact candidates are secondary findings,
not a claim that fewer inputs preserve all relevant information.

MEKD illustrates matrix-element discrimination from four-lepton
kinematics~\cite{mekd}; INFERNO optimizes learned summaries for an inference
objective~\cite{inferno}; Cowan et al.\ provide the Asimov and asymptotic
likelihood framework~\cite{cowan}. Here conventional classifier training is
retained and feature groups are compared through a common downstream fit.
No direct performance benchmark against MEKD or INFERNO has been completed.
A focused literature assessment must still establish the distinctive contribution.

\section{Data basis and suitability for publication}
The current analysis uses two simulated samples from the 2020 ATLAS Open Data
four-lepton release at $\sqrt{s}=13$ TeV: signal
\texttt{ggH125\_ZZ4lep} (DSID 345060) and background \texttt{llll} (363490).
Selection requires two electrons and two muons, trigger and lepton-quality
requirements, and $105\leq\mfl<140$ GeV. No collision data are used.

\textbf{Dataset decision.} The provider describes this collection, including
accompanying simulated samples, as educational and not suited for scientific
publications~\cite{data}. MC-only use does not establish an exception.
Should this analysis remain a preliminary study, with publication work moved
to suitable simulation, or should the proposed methodological use first be
clarified with the provider? This is a scientific-suitability issue, not an
assertion of a legal prohibition.

\begin{center}
\begin{tabular}{lrr}
\toprule
Published sample summary & Signal & Background\\
\midrule
Source MC entries & 164,716 & 554,279\\
Selected event groups (eligible development population) & 36,068 & 2,922\\
Template-role event groups & 7,221 & 590\\
Expected template yield at $10\,\mathrm{fb}^{-1}$ & 2.7683 & 2.4979\\
\bottomrule
\end{tabular}
\end{center}
Yields include role-sampling corrections and cannot be added across roles.
The expected total at $\mu=1$ is only about 5.27 events. This restricted model
omits other Higgs production modes and additional experimental backgrounds;
\texttt{llll} alone does not establish a pure $q\bar q\to ZZ$ component.
Process, reconstruction and normalization applicability need independent review.

\newpage
\section{Core method and quantities to be compared}
\subsection*{Inputs, training and weights}
\begin{center}
\begin{tabular}{ccl}
\toprule
Group & Count & Variables\\
\midrule
A & 8 & $p_{T,i},\eta_i$ for $i=1,\ldots,4$\\
B & 4 & $m_{Z_1},m_{Z_2},\Delta R_{Z_1},\Delta R_{Z_2}$\\
C & 2 & $p_{T,4\ell},\Delta\phi_{ZZ}$\\
D & 5 & $\cos\theta^*,\cos\theta_1,\cos\theta_2,
              \phi_{\mathrm{decay}},\phi_{\mathrm{production}}$\\
\bottomrule
\end{tabular}
\end{center}
Leptons are ordered by decreasing $p_T$; $Z_1$ is the same-flavour,
opposite-sign pair closest to the nominal $Z$ mass, and $Z_2$ the remaining
pair. Within each pair, $\Delta R=\sqrt{(\Delta\eta)^2+(\Delta\phi)^2}$.
In the four-lepton rest frame, $\theta^*$ compares $Z_1$ with the boosted
positive laboratory beam. Each helicity angle $\theta_{1,2}$ compares the
negative lepton with the direction opposite the other pair in its own
dilepton rest frame. Plane normals use $\mathbf{p}_{\ell^-}
\times\mathbf{p}_{\ell^+}$; the two $\phi$ angles run about $Z_1$ from
its decay normal to the other decay normal, and from the production normal
(beam crossed with $Z_1$) to its decay normal, respectively, in $[-\pi,\pi)$.

All 15 nonempty combinations use five network seeds (42--46), with common
event roles and a constant-score empty-set reference. Separate training,
validation, calibration, template and post-freeze assessment roles keep physical
groups together. The MLP has hidden widths 64, 64 and 32, trained with binary
cross entropy and AdamW. Validation AUC selects checkpoints within a maximum
200 epochs. Input dimension changes parameter count; capacity is not matched.

Training weights are $|w_i|$ divided by the training-class mean of $|w|$.
Signed physical weights are retained for yields and group-level variances.
Reported validation AUC also uses absolute physical weights and is a
checkpoint-selection diagnostic, not an independent test metric.

\subsection*{Score compression and statistical model}
Classifiers omit explicit $\mfl$, but correlated inputs may retain mass
information. One mass bin spans 105--140 GeV; each nonempty combination
has two score categories. Calibration proposes thresholds, with joint
calibration/template support checks selecting the feasible cut closest to
the median, not the cut minimizing $\W$. All events enter a category.
The reference is an inclusive count, not a resolved mass-spectrum fit.

For categories $c$ and processes $p\in\{s,b\}$, the fitted model is
\[
 {\cal L}(\mu,\gamma)=
 \prod_c\Pois(n_c\mid\mu s_c\gamma_{s,c}+b_c\gamma_{b,c})
 \prod_{p,c}\Pois(a_{p,c}\mid\tau_{p,c}\gamma_{p,c}),
 \qquad \tau_{p,c}=\frac{y_{p,c}^{\,2}}{V_{p,c}} .
\]
Here $y=\sum_g W_g$ and $V=\sum_g W_g^2$ use physical-group yield contributions.
Independent process/category Poisson auxiliary constraints model finite-template
statistics via pyhf \texttt{shapesys}; the nonnegative $\gamma$ parameters are
profiled. This is an effective-count approximation for signed MC, not a
validated complete generative model or a full experimental systematic model.

Intervals invert the profile-likelihood ratio using a one-degree-of-freedom
$\chi^2$ threshold, with $\mu\geq0$ and numerical search ceiling 20.
Failure to find an upper crossing is recorded rather than assigned a finite width.
$\W$ is the 68\% interval width
on each model's Asimov expectation at $\mu=1$; coverage is a separate question.
For each seed, $v(S)=-\W(S)$ defines exact Shapley contributions and conditional
pair interactions. These attribute the whole procedure's performance, not
intrinsic or causal physical information.

\newpage
\section{Minimum results and their interpretation}
\textbf{Observed pattern.} BC and AC are compact candidates worth further
investigation: their nominal widths are smaller than ABCD's, and the leading
configuration differs between validation AUC and interval width. The table
makes the size and limitations of those observations explicit.

\begin{center}
\textbf{Table 1. Existing exploratory results at injected $\mu=1$.}\\[0.6em]
\textit{(a) Nominal performance: median [minimum, maximum] over five seeds.}\\[0.3em]
\begin{tabular}{lcll}
\toprule
Configuration & Inputs & Validation AUC & Nominal $\W$\\
\midrule
Inclusive count & 0 & Not applicable & 1.66372 [1.66372, 1.66372]\\
ABCD & 19 & 0.84635 [0.83018, 0.85500] & 1.53593 [1.52947, 1.54805]\\
BC & 6 & 0.81314 [0.81220, 0.81425] & 1.51162 [1.51060, 1.51960]\\
AC & 10 & 0.87754 [0.87194, 0.88068] & 1.51524 [1.49720, 1.51785]\\
\bottomrule
\end{tabular}\\[1em]
\textit{(b) Paired contrasts and conditional coverage.}\\[0.3em]
\begin{tabular}{lrrc}
\toprule
Configuration & Median $\Delta\W$ & 95\% MC percentile range
 & 68\% coverage\\
\midrule
Inclusive count & --- & --- & 0.628\\
ABCD & Reference & --- & 0.646\\
BC & $-0.01995$ & [$-0.05611$, $+0.00349$] & 0.658\\
AC & $-0.02191$ & [$-0.04826$, $+0.00162$] & 0.638\\
\bottomrule
\end{tabular}
\end{center}

\textbf{How to read the table.}
$\Delta\W=\W(\mathrm{candidate})-\W(\mathrm{ABCD})$ is calculated within
each seed, then summarized by its median; negative favours the candidate.
It is not the difference between the medians in panel (a). Seed ranges
describe training variability on shared MC, not confidence intervals.

The MC ranges use all 200 valid event-group bootstrap replicas, with shared
calibration/template resampling across candidates and threshold reselection.
Networks remain fixed. These ranges do not cover retraining, assessment-parent
uncertainty or all physical variations, and are not calibrated total confidence
intervals. Both compact-versus-full ranges include zero; this establishes
neither superiority, equivalence nor noninferiority.

Coverage entries are five-seed medians from the post-freeze assessment-parent
diagnostic at $\mu=1$, with 500 pseudo-experiments per candidate per seed.
These cells have no invalid candidate fits. They are conditional checks using
the same underlying MC, not independent repetitions of the complete study.
For scale, a single independent set of 500 Bernoulli trials at coverage 0.68
has standard error $\sqrt{0.68(1-0.68)/500}\simeq0.021$.
This is not an error bar on a shared-MC seed median or on nested recalibration
trials.

\textbf{What the evidence currently supports.}
The analysis identifies procedure-dependent differences in group combinations
and a discrepancy between the leading AUC and width configurations. These
provide a concrete starting point for studying representation effects and
their reliability. BC and AC were highlighted after inspecting the rankings,
so their role here is exploratory.

\textbf{What remains unresolved.}
The finite-MC contrasts are inconclusive and the displayed coverage medians
are below 0.68. The whole threshold-selection and interval procedure needs
validation; asymptotic intervals at the small expected counts cannot be assumed
reliable. Assessment access has only a non-independent review. Pseudo-experiment
comparisons also use artificial shared-random-number coupling rather than
physical-event pairing, with allocation sensitivity still pending.
Nominal narrowing therefore does not yet establish a calibrated precision gain.
This result does not demonstrate robust compactness or a saving in training MC.

\newpage
\section{Research priorities and proposed writing direction}
\subsection*{Order of work to discuss}
\begin{enumerate}
\setlength{\itemsep}{0.15em}\setlength{\parskip}{0pt}\setlength{\parsep}{0pt}
\item \textbf{Data suitability and positioning.} Decide whether to retain this
educational sample as preliminary work, clarify methodological use with the
provider, or move publication work to suitable simulation.
\item \textbf{Interval reliability.} Independently check the signed-template
approximation and coverage of the complete selection and interval procedure.
More conditional trials cannot resolve missing sample independence.
\item \textbf{Baselines matching the claim.} A gain beyond mass-spectrum inference
requires a resolved mass baseline; otherwise retain the single-bin scope.
A physics reference needs matched inputs and validated implementation.
\item \textbf{Representation controls.} Capacity and training-size studies are
needed for claims of stable compactness or resource savings. Noninferiority
requires a prespecified acceptable loss, decision rule and independent confirmation.
\end{enumerate}
These extensions are proposed. Confirmation requires eligible simulation not
compromised by prior use; renaming or repartitioning inspected events cannot
create independence.

\subsection*{A paper organised around questions}
\textbf{Motivation:} What gap remains after existing kinematic and inference-aware
work? \textbf{Data and method:} Which population and assumptions make the comparison
appropriate? \textbf{Results and reliability:} Which groups change the outcome,
where do metric rankings differ, and which patterns survive completed checks?
\textbf{Interpretation:} What is specific to this learner, sample and likelihood,
and what evidence would support a broader claim?

The intended contribution concerns feature-group effects and their robustness.
Compact candidates motivate follow-up, not the main conclusion. Novelty remains
to be established, and an inconclusive result is possible.

\subsection*{Guidance requested}
\begin{enumerate}
\setlength{\itemsep}{0.15em}\setlength{\parskip}{0pt}\setlength{\parsep}{0pt}
\item \textbf{Research design:} Is this question worthwhile and sufficiently
focused? What data, baseline and methodological changes are essential?
\item \textbf{Objectives and conclusions:} What should be the primary estimand
and minimum evidence for a defensible conclusion? Which validation comes first?
\item \textbf{Manuscript development:} How should the contribution be framed,
which sections deserve emphasis, and how should the present conclusions be worded?
\end{enumerate}

\begingroup
\small
\begin{thebibliography}{4}
\setlength{\itemsep}{0pt}
\setlength{\parskip}{0pt}
\bibitem{data} ATLAS Collaboration, 2020 four-lepton collection,
\href{https://opendata.cern.ch/record/15005}{CERN Open Data record 15005}
(accessed 4 October 2026).
\bibitem{mekd} P.~Avery et al., \textit{Precision studies of the Higgs boson
decay channel $H\to ZZ\to4\ell$ with MEKD}, \arxivref{1210.0896}.
\bibitem{inferno} P.~de Castro and T.~Dorigo, \textit{INFERNO:
Inference-Aware Neural Optimisation}, \arxivref{1806.04743}.
\bibitem{cowan} G.~Cowan et al., \textit{Asymptotic formulae for likelihood-based
tests of new physics}, \arxivref{1007.1727}.
\end{thebibliography}
\textit{Local evidence:} the selected test05 aggregate snapshot and
\texttt{paper/result-evidence.md}; no new experiment was run for this outline.
Time and resources will be assessed after the scope is agreed.
\endgroup
\end{document}
